\begin{table}[t]
\centering
\caption{Standardized Distributional Effect (SDE) Appendix}
\label{tab:sde}
\begin{adjustbox}{max width=\textwidth}
\begin{tabular}{lcccccc}
\hline\hline
Outcome & $\hat{\beta}$ & SE & SD(Y) & SDE & SE(SDE) & Classification \\
\hline
\textit{Panel A: Pooled} & & & & & & \\
Pro-Worker Verdict & -0.0601 & 0.0121 & 0.419 & -0.1435 & 0.0289 & Moderate \\
Full Proced\^{e}ncia & 0.0144 & 0.0165 & 0.338 & 0.0427 & 0.0489 & Small \\
{}[0.5em]
\textit{Panel B: Heterogeneous} & & & & & & \\
High Discretion & -0.0485 & 0.0142 & 0.400 & -0.1213 & 0.0355 & Moderate \\
Low Discretion & -0.0777 & 0.0194 & 0.380 & -0.2044 & 0.0510 & Large \\
\hline\hline
\end{tabular}
\end{adjustbox}
\begin{flushleft}\small
\textbf{Country:} Brazil.
\textbf{Research question:} Does the 2017 labor reform compress cross-court heterogeneity in adjudication?
\textbf{Policy mechanism:} Lei 13.467 shifted litigation costs to losing plaintiffs, changing the equilibrium filing pool and potentially disciplining judicial discretion.
\textbf{Outcome definition:} Pro-worker verdict (Proced\^{e}ncia or Proced\^{e}ncia em Parte vs.\ Improced\^{e}ncia).
\textbf{Treatment:} Pre-reform vara leniency $L_j$ (empirical Bayes-shrunk, standardized) interacted with post-reform indicator.
\textbf{Data:} DataJud API (CNJ), TRT2/TRT4/TRT15, 2012--2023.
\textbf{Method:} Case-level OLS with vara and year-month fixed effects; standard errors clustered by vara.
\textbf{Sample:} First-instance labor cases assigned by lottery (\textit{sorteio}).
Classification refers to magnitude, not statistical significance. Large: $|\text{SDE}| > 0.15$; Moderate: $0.05$--$0.15$; Small: $0.005$--$0.05$; Null: $< 0.005$.
\end{flushleft}
\end{table}
